\documentclass{report}
\usepackage[utf8]{inputenc}
\usepackage[spanish, es-noshorthands]{babel}
\usepackage{amsmath}
\usepackage{amssymb}
\usepackage{amsthm}
\usepackage{tikz}
\usepackage[colorlinks=true, allcolors=blue]{hyperref}
\usetikzlibrary{automata, positioning}
\interfootnotelinepenalty=10000

\theoremstyle{definition}
\newtheorem{definicion}{Definición}[chapter]
\theoremstyle{plain}
\newtheorem{teorema}{Teorema}[chapter]
\newtheorem{lema}{Lema}[chapter]
\theoremstyle{remark}
\newtheorem{ejemplo}{Ejemplo}[chapter]
\input{../word-comments.tex}

\begin{document}

\setcounter{chapter}{1}
\chapter{Gramáticas de concatenación de rango y el SAT como lenguaje formal}
\label{cap:rcg-sat}

Las gramáticas de concatenación de rango son el formalismo \quitaesto{de base}
sobre el que se desarrolla el resto del trabajo. \quitaesto{Los conceptos básicos
  que se utilizan en este capítulo se presentaron en el capítulo
  anterior.} \agregaesto{MISMO PARRAFO}El capítulo se organiza en tres secciones. La sección~\ref{sec:rcg}
presenta las gramáticas de concatenación de rango, su \comment{presentación}{presenta y presentación muy seguidas una de otra}
intuitiva, sus definiciones formales y un ejemplo \quitaesto{completo} de
reconocimiento. La sección~\ref{sec:propiedades-rcg} expone las
propiedades de clausura del formalismo y su caracterización en
términos de complejidad. La sección~\ref{sec:antecedentes} revisa los
antecedentes directos del trabajo: los estudios previos \quitaesto{en} \agregaesto{de} la Facultad \agregaesto{de Matemática y Coputación de La Universidad de La Habana}
que abordan el SAT mediante formalismos de la teoría de lenguajes.

\section{Gramáticas de concatenación de rango}
\label{sec:rcg}

Las gramáticas de concatenación de rango (RCG)~\cite{boullier1998proposal}
son un formalismo propuesto en 1998 por Pierre Boullier, un investigador
en el campo de la lingüística computacional. El objetivo inicial del
formalismo era proporcionar un modelo más expresivo que las gramáticas
libres del contexto para describir ciertos fenómenos del lenguaje
natural, como los números chinos y el orden variable de algunas palabras
en alemán~\cite{boullier1999chinese}.

Antes de entrar a la presentación del formalismo, es necesario precisar
una convención terminológica que rige a lo largo de todo el documento.
El formalismo que se presenta en este capítulo corresponde, en la
literatura, a las \textbf{gramáticas de concatenación de rango
  positivas} (pRCG), una de las dos variantes del formalismo general de
Boullier. La otra variante, las gramáticas de concatenación de rango con
negación, admite cláusulas con predicados negados y no se utiliza en
este trabajo. La restricción a la variante positiva es importante
porque, como se establece más adelante, la clase de lenguajes
reconocidos por una pRCG coincide exactamente con la clase $P$, y
los resultados de los capítulos posteriores se apoyan en dicha
caracterización. Para abreviar, a lo largo del documento se adopta el
nombre \textbf{gramática de concatenación de rango} (RCG) para
referirse, salvo mención explícita, a la variante positiva.

\subsection{Presentación intuitiva}

\comment{}{no lo se rick, presentación intuitiva no me gusta, mejor pon `idea intuitiva' o algo así} A continuación se introducen las nociones de predicado, argumento,
sustitución de rango y derivación, que se utilizan más adelante.

A los no terminales de una RCG se les llama \textbf{predicados}. Cada
predicado tiene un conjunto de argumentos. A la cantidad de argumentos
de un predicado se le denomina \textbf{aridad}. Por ejemplo, $A(X, Y)$
denota \comment{el predicado $A$ con argumentos $X$ e $Y$; en este caso, la
aridad de $A$ es $2$}{escribelo mejor como `el predicado a de aridad 2 con argumentos $X$ e $Y$'}.

Cada argumento de un predicado puede estar formado por variables y
terminales. Por convenio, las variables se denotan por letras mayúsculas
del final del alfabeto, y los terminales con letras minúsculas. \comment{Con este
convenio, el predicado $B(aX, XY, abZ)$ tiene aridad $3$}{en este ejemplo y en el anterior pusiste predicados con $n$ argumentos y $n$ variables, ahora, de la forma que lo escribes no se si es casualidad, o es algo que exige el formalismo, no queda claro. Si es posible tener mas o menos argumentos que variables, deberías buscarte un ejemplo de ello, de lo contrario, creo que deberías decir explicitamente que debe ser la misma cantidad}. El primer
argumento está formado por el terminal $a$ y la variable $X$. El
segundo argumento está formado por la concatenación de las variables
$X$ e $Y$. El tercer argumento está formado por la concatenación de los
terminales $a$ y $b$ y la variable $Z$.

Cada predicado reconoce un vector de cadenas cuya dimensión es la
aridad del predicado. Cada cadena del vector se asocia a un argumento
del predicado.

Por ejemplo, si al predicado $A(X, Y)$ se le asocia el vector $[abc,
      d]$, al primer argumento de $A$ se le asigna la cadena $abc$ y al
segundo la cadena $d$. Esto significa que $X = abc$ y que $Y = d$.

De manera similar, si al predicado $B(aX, XY, abZ)$ se le asigna el
vector $[aa, de, abcc]$, al primer argumento de $B$ se le asigna la
cadena $aa$, al segundo la cadena $de$ y al tercero la cadena $abcc$.
Es decir, $aX = aa$, $XY = de$ y $abZ = abcc$. Como $X$ es una variable
y se tiene que $aX = aa$, el único valor que puede tomar $X$ es $X = a$.
En el tercer argumento ocurre algo similar: como $abZ = abcc$, para que
se cumpla la igualdad el único valor posible para $Z$ es $Z = cc$. En el
segundo argumento, las variables $X$ e $Y$ deben tomar valores tales que
su concatenación sea $de$; esto se puede hacer de tres formas:
$X = d$, $Y = e$; $X = de$, $Y = \varepsilon$; y $X = \varepsilon$,
$Y = de$.

\comment{Si al predicado $B(aX, XY, abZ)$ se le asigna el vector
$[bb, de, abcc]$, entonces al tenerse $aX = bb$, $X$ no puede tomar
ningún valor: el primer carácter $b$ de la cadena no coincide con el
terminal $a$ que encabeza el argumento $aX$. Por tanto, el predicado
$B$ no reconoce este vector.}{este párrafo deberías escribirlo al reves, en plan... `el predicado $B$ no reconoce el vector tal, porque tal cosa', no digas que `se le asigna el vector' si todavia al final no se reconoce, tal parece que vas a explicar como se reconoce ese vector y al final resulta que no se pudo.}

A las producciones de esta gramática se les denomina \textbf{cláusulas}.
La parte izquierda de la producción siempre está formada por un único
predicado. La parte derecha puede contener cualquier cantidad de
predicados con sus respectivos argumentos, o la cadena vacía. \comment{Un ejemplo
de cláusula es}{ah mira, pones el ejemplo aquí, pero como diría Fernando, \lquoti la gente no lee con \textbf{lookahead}\rquoti XD}
\[
  A(XYZ, W) \to B(X)\, C(XY, Z)\, D(W).
\]

La regla anterior tiene el siguiente significado: el predicado $A$
recibe un vector de dimensión $2$; a partir de las cadenas del vector,
se asignan valores a las variables $X$, $Y$, $Z$ y $W$, y con esos
valores se construyen los vectores que reciben los predicados $B$, $C$
y $D$.

El proceso se ilustra con un ejemplo en el que el predicado $A$ recibe
el vector de cadenas $[abc, d]$. El primer paso es asignar las cadenas
del vector a los argumentos del predicado. El primer argumento de $A$
es $abc$ y el segundo es $d$. Como el primer argumento de $A$ está
definido como $XYZ$ y el segundo como $W$, debe cumplirse que $XYZ = abc$
y $W = d$. Esto significa que las variables $X$, $Y$ y $Z$ deben tomar
los valores posibles cuya concatenación sea $abc$, mientras que $W$ solo
puede tomar el valor $d$. En el cuadro~\ref{tab:valores-xyzw} se
muestran algunas asignaciones posibles para las variables $X$, $Y$, $Z$
y $W$ a partir del vector de entrada.

\begin{table}[ht]
  \centering
  \begin{tabular}{|c|c|c|c|}
    \hline
    $X$           & $Y$           & $Z$           & $W$ \\ \hline
    $\varepsilon$ & $\varepsilon$ & $abc$         & $d$ \\ \hline
    $\varepsilon$ & $a$           & $bc$          & $d$ \\ \hline
    $\varepsilon$ & $ab$          & $c$           & $d$ \\ \hline
    $\varepsilon$ & $abc$         & $\varepsilon$ & $d$ \\ \hline
    $a$           & $\varepsilon$ & $bc$          & $d$ \\ \hline
    $a$           & $b$           & $c$           & $d$ \\ \hline
    $a$           & $bc$          & $\varepsilon$ & $d$ \\ \hline
    $ab$          & $\varepsilon$ & $c$           & $d$ \\ \hline
    $ab$          & $c$           & $\varepsilon$ & $d$ \\ \hline
    $abc$         & $\varepsilon$ & $\varepsilon$ & $d$ \\ \hline
  \end{tabular}
  \caption{Asignaciones posibles para las variables $X$, $Y$, $Z$ y
    $W$ a partir del vector $[abc, d]$.}
  \label{tab:valores-xyzw}
\end{table}

\comment{Para continuar con el ejemplo}{aquí creo que queda mejor `Continuando con el ejemplo'}, se supone que los valores de $X$, $Y$ y
$Z$ son $X = ab$, $Y = \varepsilon$ y $Z = c$. A partir de esta
asignación se construyen los vectores con los que se evalúan los
predicados del lado derecho $B(X)\, C(XY, Z)\, D(W)$. Como $X = ab$,
$Y = \varepsilon$, $Z = c$ y $W = d$, el lado derecho se evalúa con los
vectores $[ab]$, $[ab, c]$ y $[d]$, y el resultado es la derivación
$B(ab)\, C(ab, c)\, D(d)$.

El proceso de asignar los vectores a los argumentos, identificar los
valores de las variables e instanciar las partes derechas se repite en
cada uno de los predicados del lado derecho de la cláusula.

Existe otro tipo de producciones, denominadas cláusulas terminales, en
las que un predicado deriva en $\varepsilon$ para determinados valores
de sus argumentos. Por ejemplo,
\[
  B(ab) \to \varepsilon, \qquad
  C(ab, c) \to \varepsilon, \qquad
  D(d) \to \varepsilon.
\]
\comment{La forma general de estas producciones es
$A(X_1, \dots, X_n) \to \varepsilon$.}{me choca un poco que en los ejemplos pusiste $B(ab)$ y abajo pones que la forma general es esta, con las $X$ grandes, que cosa son esas $X$? cadenas? variables? no queda claro}

Un predicado reconoce un vector de cadenas si existe una asignación de
sus argumentos a valores de las variables y una secuencia de
derivaciones que termina en la cadena vacía. Por ejemplo, en el caso de
$B(ab) \to \varepsilon$, el predicado $B$ reconoce la cadena $ab$.

A modo de ilustración, se considera el siguiente conjunto de
producciones:
\begin{enumerate}
  \item $A(XYZ, W) \to B(X)\, C(XY, Z)\, D(W)$,
  \item $B(ab) \to \varepsilon$,
  \item $C(ab, c) \to \varepsilon$,
  \item $D(d) \to \varepsilon$.
\end{enumerate}
Con estas producciones, el predicado $A$ reconoce el vector $[abc, d]$:
con la asignación $X = ab$, $Y = \varepsilon$, $Z = c$ y $W = d$, el
predicado $A$ deriva en $B(ab)\, C(ab, c)\, D(d)$, y cada uno de los
predicados $B(ab)$, $C(ab, c)$ y $D(d)$ deriva en la cadena vacía.

Un predicado no reconoce un vector de cadenas cuando para ninguna
asignación de variables se deriva en la cadena vacía.

Presentadas estas nociones, a continuación se introducen las
definiciones formales de las gramáticas de concatenación de rango.

\subsection{Definiciones formales}

\comment{}{aquí sería bueno poner aunque sea una oración introductoria para no entrar así tan áspero}

\begin{definicion}
  \comment{Un \textbf{rango} es una tupla $(i, j)$ que representa un intervalo de
  posiciones en una cadena, donde $i$ y $j$ son enteros no negativos tales
  que $i \le j$.}{esto es algo muy personal mío, no entiendo para que defines lo que es un rango? abajo dices sustitución de rango pero sustituyes por una subcadena. De la forma que está escrito este concepto, tal parece que las sustituciones de rango sustituyen una variable por una tupla, lo que quieres decir es que lo sustituyes por la cadena de $i$ a $j$. Yo revisaría como presentas esto.}
\end{definicion}

Por ejemplo, con índices indexados en $0$, para la cadena $abcd$ el
rango $(1, 3)$ representa la subcadena $bc$.

\begin{definicion}
  \label{def:rcg}
  Una \textbf{gramática de concatenación de rango} se define como una
  $5$-tupla
  \[
    G = (N, T, V, P, S),
  \]
  donde:
  \begin{itemize}
    \item $N$ es un conjunto finito de \textbf{predicados} o símbolos no
          terminales; cada predicado tiene una aridad, que indica la dimensión
          del vector de cadenas que reconoce;
    \item $T$ es un conjunto finito de \textbf{símbolos terminales};
    \item $V$ es un conjunto finito de \textbf{variables};
    \item $P$ es un conjunto finito de \textbf{cláusulas} de la forma
          \[
            A(x_1, x_2, \dots, x_k) \to B_1(y_{1,1}, y_{1,2}, \dots, y_{1,m_1})
            \dots B_n(y_{n,1}, y_{n,2}, \dots, y_{n,m_n}),
          \]
          donde $A, B_i \in N$, $x_i, y_{i,j} \in (V \cup T)^*$, $n \ge 0$,
          $k$ es la aridad de $A$ y $m_i$ es la aridad de $B_i$; cuando
          $n = 0$, la parte derecha es la cadena vacía y la cláusula se
          denomina terminal;
    \item $S \in N$ es el \textbf{predicado inicial} de la gramática, que
          siempre tiene aridad $1$.
  \end{itemize}
\end{definicion}

Por ejemplo, a continuación se muestra una gramática de concatenación
de rango que reconoce el lenguaje $L_{\mathrm{copy}}^3 = \{ w^3 \mid w
  \in \{a, b, c\}^* \}$:
\[
  G_{\mathrm{copy}}^3 = (N, T, V, P, S),
\]
donde $N = \{A, S\}$, $T = \{a, b, c\}$, $V = \{X, Y, Z\}$, el símbolo
inicial es $S$ y el conjunto de cláusulas $P$ es el siguiente:
\begin{enumerate}
  \item $S(XYZ) \to A(X, Y, Z)$,
  \item $A(aX, aY, aZ) \to A(X, Y, Z)$,
  \item $A(bX, bY, bZ) \to A(X, Y, Z)$,
  \item $A(cX, cY, cZ) \to A(X, Y, Z)$,
  \item $A(\varepsilon, \varepsilon, \varepsilon) \to \varepsilon$.
\end{enumerate}

A diferencia de las gramáticas presentadas en el capítulo anterior, que
se definen a partir de un mecanismo de generación de cadenas, las RCG
se formulan directamente como un mecanismo de reconocimiento. Su
funcionamiento consiste en determinar si una cadena pertenece o no al
lenguaje.

\begin{definicion}
  Una \textbf{sustitución de rango} es un mecanismo que reemplaza una
  variable por un rango de la cadena, respetando la estructura del
  argumento que se asocia a la cadena que se reconoce.
\end{definicion}

Por ejemplo, dado el predicado $A(Xa)$ con $X \in V$ y $a \in T$, la
estructura del argumento de $A$ es una variable $X$ seguida del
terminal $a$. Si el predicado $A$ recibe la cadena $baa$, entonces $X$
se puede asociar con la subcadena $ba$ de la cadena original: si
$X = ba$, entonces $Xa = baa$. La variable $X$ no puede tomar el valor
$baa$, porque ningún carácter de la cadena de entrada coincide con el
terminal $a$ al final del argumento. La variable $X$ tampoco puede
tomar el valor $b$, porque con ese valor el argumento $Xa$ no cubre la
cadena completa.

\begin{definicion}
  Las \textbf{gramáticas de concatenación de rango simple} (sRCG) son un
  subconjunto de las RCG que restringen la forma de las cláusulas. Una
  RCG $G$ es \textbf{simple} si los argumentos del lado derecho de cada
  cláusula son variables distintas, y todas estas variables aparecen una
  sola vez en los argumentos del lado izquierdo.
\end{definicion}

Por ejemplo, la cláusula $A(XY) \to B(X)\, C(Y)$ cumple la condición de
simplicidad: los argumentos del lado derecho, $X$ e $Y$, son variables
distintas, y cada una aparece una sola vez en el lado izquierdo. La
cláusula $A(X) \to B(X)\, C(X)$ no la cumple, porque la variable $X$
aparece en dos argumentos del lado derecho; una RCG que la contenga no
es simple.

Las sRCG son equivalentes a los sistemas lineales de reescritura libres
del contexto (LCFRS) introducidos en~\cite{vijayshanker1987lcfrs}. Este
subconjunto se utiliza en~\cite{aguileralopez2016} para una variante
del SAT denominada \textbf{Satisfacibilidad Booleana de Concatenación
  de Rango Simple}, que se discute en la sección~\ref{sec:antecedentes}.

Una vez introducidas las definiciones formales, a continuación se
presenta un ejemplo completo de reconocimiento de una cadena por una
RCG, aplicando los mecanismos descritos en esta sección.

\subsection{Ejemplo de reconocimiento}

La presentación intuitiva describió cómo se asocian los elementos de
un vector a los argumentos de un predicado y cómo se construyen los
vectores que reciben los predicados del lado derecho de una cláusula.
En esta subsección se aplica ese proceso a un caso completo: el
reconocimiento de la cadena $abcabcabc$ por la gramática
$G_{\mathrm{copy}}^3$ presentada anteriormente.

La cadena $abcabcabc$ se reconoce por $G_{\mathrm{copy}}^3$ mediante
la siguiente secuencia de derivaciones:
\[
  S(abcabcabc) \to A(abc, abc, abc) \to A(bc, bc, bc) \to A(c, c, c)
  \to A(\varepsilon, \varepsilon, \varepsilon) \to \varepsilon.
\]
A continuación se describen los pasos.
\begin{itemize}
  \item \textbf{Primer paso:} en la primera cláusula
        $S(XYZ) \to A(X, Y, Z)$, la sustitución de rango asocia las
        variables $X$, $Y$, $Z$ a los valores $X = abc$, $Y = abc$ y
        $Z = abc$. De esta forma se deriva en el predicado $A(abc, abc,
          abc)$.
  \item \textbf{Segundo paso:} en la segunda cláusula
        $A(aX, aY, aZ) \to A(X, Y, Z)$, la sustitución de rango asocia las
        variables $X$, $Y$, $Z$ a los valores $X = bc$, $Y = bc$ y
        $Z = bc$. Con estos valores se deriva en el predicado
        $A(bc, bc, bc)$.
  \item \textbf{Tercer paso:} en la tercera cláusula
        $A(bX, bY, bZ) \to A(X, Y, Z)$, la sustitución de rango asocia las
        variables $X$, $Y$, $Z$ a los valores $X = c$, $Y = c$ y $Z = c$.
        Con estos valores se deriva en el predicado $A(c, c, c)$.
  \item \textbf{Cuarto paso:} en la cuarta cláusula
        $A(cX, cY, cZ) \to A(X, Y, Z)$, la sustitución de rango asocia las
        variables $X$, $Y$, $Z$ a los valores $X = \varepsilon$,
        $Y = \varepsilon$ y $Z = \varepsilon$. Con estos valores se deriva
        en el predicado $A(\varepsilon, \varepsilon, \varepsilon)$.
  \item \textbf{Quinto paso:} en la última cláusula
        $A(\varepsilon, \varepsilon, \varepsilon) \to \varepsilon$ se deriva
        en la cadena vacía, con lo que la cadena $abcabcabc$ queda
        reconocida.
\end{itemize}

Una vez presentado el proceso de derivación, en la siguiente sección se
analizan las propiedades del formalismo: clausura bajo intersección,
comparación con las gramáticas libres del contexto y complejidad del
problema de la palabra.

\section{Propiedades de las RCG}
\label{sec:propiedades-rcg}

Las RCG son cerradas bajo unión, intersección y
complemento~\cite{boullier1998proposal}. De estas tres operaciones, la
clausura bajo intersección es la que importa para este trabajo. En el
enfoque que
se revisa en la sección~\ref{sec:antecedentes}, la satisfacibilidad de
una fórmula se decide según si una intersección de dos lenguajes es
vacía. En esta sección se presentan estas propiedades de clausura y la
caracterización del reconocimiento en términos de complejidad. La
clausura bajo intersección se usa en el capítulo~\ref{cap:interseccion},
donde se construye una RCG cuyo lenguaje es esa intersección. El
reconocimiento polinomial se usa en el
capítulo~\ref{cap:vacio-interseccion-rcg-regular}, donde se caracteriza la
decisión de su vacío.

\begin{teorema}[Clausura bajo intersección]
  \label{teo:clausura-rcg}
  Si $L_1$ y $L_2$ son lenguajes reconocidos por RCG, entonces
  $L_1 \cap L_2$ también lo es~\cite{boullier1998proposal}.
\end{teorema}

Las RCG son más expresivas que las gramáticas libres del contexto. Por
ejemplo, para todo $k \ge 2$, el lenguaje $L_{\mathrm{copy}}^k = \{ w^k
  \mid w \in \Sigma^*\}$ es dependiente del contexto y no libre del
contexto~\cite{hopcroft2006introduction}, y se reconoce por una RCG. En
particular, el lenguaje $L_{\mathrm{copy}}^3$ se reconoce por la
gramática $G_{\mathrm{copy}}^3$ del ejemplo anterior.

\comment{Una vez presentadas la clausura bajo intersección y la comparación con
las gramáticas libres del contexto,}{no sigas poniendo las transiciones así, yo no se decirte exactamente que regla gramatical estás violando ahí, pero se siente discordante la oración. Es como si la pusieras en futuro y en presente al mismo tiempo, me suena rara} en la próxima subsección se analiza
el problema de la palabra de las RCG.

\subsection{Complejidad y caracterización en términos de $P$}

El siguiente resultado es una caracterización del reconocimiento de las
RCG en términos de complejidad.

\begin{teorema}[Caracterización en términos de $P$]
  \label{teo:prcg-ptime}
  \comment{Un lenguaje $L$ se reconoce por una RCG en su variante positiva si y
  solo si $L$ admite un algoritmo de reconocimiento que se ejecuta en
  tiempo polinomial~\cite{bertsch2001rcg, boullier2000range}.}{revisa este teorema. De la forma en que lo escribes quedan las cosas quedan como que no tienen que ver. O sea, me dices que para que se reconozca un lenguaje, tiene que pasar que L se reconozca en tiempo polinomial?...

  Creo que ya entendí lo que dices, si mal no entendí $L$ es un lenguaje que está en P. Lo que quieres decir es que $\mathcal{L}(pRGC) = P$. Igual revisa como escribes eso, que no queda claro a la primera. Para empezar, que es eso de que los problemas sean lenguajes? creo que eso no lo has definido}
\end{teorema}

La dirección que afirma
que todo lenguaje RCG está en $P$ se sigue de la existencia de un
algoritmo de reconocimiento polinomial para el formalismo, que se
describe a continuación.

El algoritmo de reconocimiento estándar para las RCG utiliza un proceso
de memorización sobre las cadenas asignadas a los argumentos de los
predicados~\cite{boullier1998proposal}. La cantidad máxima de rangos de
una cadena de longitud $n$ es $O(n^2)$, y la máxima aridad de un predicado
es constante, por lo que el proceso de memorización emplea una cantidad
polinomial de estados. La complejidad del algoritmo es $O(|P|
  n^{2h(l+1)})$, donde $h$ es la máxima aridad de un predicado, $l$ es
la máxima cantidad de predicados en el lado derecho de una cláusula y
$n$ es la longitud de la cadena que se reconoce. El teorema establece,
además, la dirección inversa: \comment{todo lenguaje en $P$}{lo mismo aquí. Cuando revises la parte de clases de problemas del primer capítulo, tienes que hacerlo de manera tal que te quede que las clases P, NP y NP-completo sean familias de lenguajes. Eso o cambiar estos párrafos} admite una RCG
que lo reconoce.

El alcance exacto
de la caracterización y sus consecuencias para la relación entre
formalismos gramaticales y el SAT son el objeto
central de los capítulos posteriores.

Una vez establecida la caracterización de las RCG como formalismo que
reconoce exactamente la clase $P$, en la siguiente sección se
revisan los antecedentes del trabajo.

\section{Antecedentes en el estudio del SAT mediante formalismos de la teoría de lenguajes}
\label{sec:antecedentes}

El estudio del problema de la satisfacibilidad booleana mediante
formalismos de la teoría de lenguajes cuenta con antecedentes en la
Facultad de Matemática y Computación de la Universidad de La Habana. El
presente trabajo se apoya en dos de ellos: el problema de la
satisfacibilidad booleana libre del contexto~\cite{fernandezarias2007},
abreviado SATCF, y el de la satisfacibilidad booleana de concatenación
de rango simple~\cite{aguileralopez2016}, abreviado SATSRC.

SATCF y SATSRC comparten una misma estrategia. A una fórmula booleana se
le asocia un lenguaje regular que representa sus interpretaciones
verdaderas, bajo el supuesto de que cada ocurrencia de variable se trata
por separado. Sobre ese lenguaje se impone después que las ocurrencias
de una misma variable tomen el mismo valor, mediante una gramática. La
satisfacibilidad se decide según si la intersección de ambos lenguajes
es vacía.

La sección se organiza según esa estrategia. En la
subsección~\ref{subsec:automata-interpretaciones} se describe el
autómata con que se construye el lenguaje regular de las
interpretaciones. En la subsección~\ref{subsec:consistencia-clase} se
presentan la gramática que impone la consistencia, la clase de fórmulas
que cada trabajo resuelve y la razón por la que las RCG son el
siguiente formalismo de esta línea de trabajo.

\subsection{El autómata de las interpretaciones}
\label{subsec:automata-interpretaciones}

El lenguaje regular de las interpretaciones se construye mediante un
autómata finito, el \textbf{autómata
  booleano}~\cite{fernandezarias2007}. Su alfabeto de entrada es
$\{0,1\}$, cuyos símbolos son los valores de verdad: $1$ el verdadero y
$0$ el falso. Cada ocurrencia de variable de la fórmula aporta uno de
esos símbolos, según su valor. Una cadena de ceros y unos describe así
una asignación a las ocurrencias, tratadas por separado, y se acepta
cuando esa asignación hace verdadera la fórmula.

El autómata de una sola ocurrencia consta de un estado inicial $q_0$ y
dos estados más, $V$ y $F$. Al leer $1$ se pasa a $V$ y al leer $0$ se
pasa a $F$. El estado $V$ es de aceptación y representa que la fórmula
resulta verdadera; el estado $F$ representa lo contrario. La
figura~\ref{fig:automata-base} lo muestra.

\begin{figure}[ht]
  \centering
  \begin{tikzpicture}[->, >=stealth, shorten >=1pt, auto, semithick]
    \node[state, initial] (q0) at (0,0) {$q_0$};
    \node[state, accepting] (V) at (2.6,1) {$V$};
    \node[state] (F) at (2.6,-1) {$F$};
    \path (q0) edge node {$1$} (V)
    (q0) edge node[swap] {$0$} (F);
  \end{tikzpicture}
  \caption{Autómata booleano de una ocurrencia de variable.}
  \label{fig:automata-base}
\end{figure}

Los autómatas de las ocurrencias se combinan según los operadores de la
fórmula~\cite{fernandezarias2007}. El autómata de una fórmula compuesta
se obtiene de los autómatas de sus subfórmulas, según el operador
principal.

La negación se aplica a una subfórmula. Sea $A$ una subfórmula y $B$ su
autómata booleano; con la negación se intercambian el estado de
aceptación $V$ y el de rechazo $F$, de modo que el autómata de $\neg A$
es $B$ con $V$ como no aceptante y $F$ como aceptante. Cada cadena que
hace verdadera a $A$ hace falsa a $\neg A$, y al revés. La
figura~\ref{fig:automata-negacion} lo muestra.

\begin{figure}[ht]
  \centering
  \begin{tabular}{cc}
    \begin{tikzpicture}[->, >=stealth, shorten >=1pt, auto, semithick]
      \node[state, initial] (q0) at (0,0) {$q_0$};
      \node[draw, rounded corners, minimum width=0.8cm, minimum height=0.9cm] (B) at (1.7,0) {$B$};
      \node[state, accepting] (V) at (3.5,0.8) {$V$};
      \node[state] (F) at (3.5,-0.8) {$F$};
      \draw (q0) -- (B);
      \draw (B) -- (V);
      \draw (B) -- (F);
    \end{tikzpicture}
                                 &
    \begin{tikzpicture}[->, >=stealth, shorten >=1pt, auto, semithick]
      \node[state, initial] (q0) at (0,0) {$q_0$};
      \node[draw, rounded corners, minimum width=0.8cm, minimum height=0.9cm] (B) at (1.7,0) {$B$};
      \node[state] (V) at (3.5,0.8) {$V$};
      \node[state, accepting] (F) at (3.5,-0.8) {$F$};
      \draw (q0) -- (B);
      \draw (B) -- (V);
      \draw (B) -- (F);
    \end{tikzpicture}
    \\[0.4em]
    {\small (a) autómata de $A$} & {\small (b) autómata de $\neg A$}
  \end{tabular}
  \caption{La negación intercambia el papel de $V$ y $F$: en $\neg A$ el
    estado aceptante es $F$.}
  \label{fig:automata-negacion}
\end{figure}

La conjunción y la disyunción se aplican a dos subfórmulas. Sean $A_1$ y
$A_2$ dos subfórmulas y $B_1$, $B_2$ sus autómatas booleanos, con estados
de aceptación $V_1$, $V_2$ y de rechazo $F_1$, $F_2$.

En la disyunción $A_1 \vee A_2$ se concatenan los autómatas de $A_1$ y de
$A_2$. Si en el autómata de $A_1$ se alcanza el estado $V_1$, la
disyunción ya resulta verdadera; desde ahí, a lo largo de un \textbf{camino
  de extensión}\footnote{Una sucesión de estados en la que se consumen los
  símbolos restantes de la cadena y se llega al estado ya determinado,
  cualquiera que sea su valor.}, se consumen los símbolos de $A_2$ hasta el
estado $V$ del autómata combinado. Si se alcanza el estado $F_1$, el valor depende de
$A_2$, y por eso el estado $F_1$ se une con el estado inicial del
autómata de $A_2$; los estados $V_2$ y $F_2$ pasan a ser los del autómata
combinado. La figura~\ref{fig:automata-disyuncion} lo muestra.

\begin{figure}[ht]
  \centering
  \begin{tikzpicture}[->, >=stealth, shorten >=1pt, auto, semithick]
    \node[state, initial] (q0) at (0,0) {$q_0$};
    \node[draw, rounded corners, minimum width=0.8cm, minimum height=0.9cm] (B1) at (1.7,0) {$B_1$};
    \node[state] (V1) at (3.4,1.1) {$V_1$};
    \node[state] (F1) at (3.4,-1.1) {$F_1$};
    \node[draw, rounded corners, minimum width=0.8cm, minimum height=0.9cm] (B2) at (5.1,-1.1) {$B_2$};
    \node[state, accepting] (V) at (7.6,1.1) {$V$};
    \node[state] (F) at (7.6,-1.1) {$F$};
    \draw (q0) -- (B1);
    \draw (B1) -- (V1);
    \draw (B1) -- (F1);
    \draw[dashed] (V1) -- (V);
    \draw (F1) -- (B2);
    \draw (B2) -- (V);
    \draw (B2) -- (F);
  \end{tikzpicture}
  \caption{Disyunción $A_1 \vee A_2$. Al alcanzar $V_1$, un camino de
    extensión (flecha discontinua) lleva al estado aceptante $V$; al
    alcanzar $F_1$, el cómputo continúa en $B_2$.}
  \label{fig:automata-disyuncion}
\end{figure}

La conjunción $A_1 \wedge A_2$ es \queesesto{simétrica}. Si en el autómata de $A_1$
se alcanza el estado $F_1$, la conjunción ya resulta falsa, y a lo largo
de un camino de extensión se consumen los símbolos de $A_2$ hasta el
estado $F$ del autómata combinado. Si se alcanza el estado $V_1$, el
valor depende de $A_2$, y el estado $V_1$ se une con el estado inicial
del autómata de $A_2$. \comment{La figura~\ref{fig:automata-conjuncion} lo
muestra.}{esto es algo muy personal mío tambien, pero no me gusta como referencias las figuras. Mejor ponlas al principio del párrafo: en la figura nosecual se ve que la conjunción es noseque. blablabla}

\begin{figure}[ht]
  \centering
  \begin{tikzpicture}[->, >=stealth, shorten >=1pt, auto, semithick]
    \node[state, initial] (q0) at (0,0) {$q_0$};
    \node[draw, rounded corners, minimum width=0.8cm, minimum height=0.9cm] (B1) at (1.7,0) {$B_1$};
    \node[state] (V1) at (3.4,1.1) {$V_1$};
    \node[state] (F1) at (3.4,-1.1) {$F_1$};
    \node[draw, rounded corners, minimum width=0.8cm, minimum height=0.9cm] (B2) at (5.1,1.1) {$B_2$};
    \node[state, accepting] (V) at (7.6,1.1) {$V$};
    \node[state] (F) at (7.6,-1.1) {$F$};
    \draw (q0) -- (B1);
    \draw (B1) -- (V1);
    \draw (B1) -- (F1);
    \draw (V1) -- (B2);
    \draw[dashed] (F1) -- (F);
    \draw (B2) -- (V);
    \draw (B2) -- (F);
  \end{tikzpicture}
  \caption{Conjunción $A_1 \wedge A_2$. Al alcanzar $F_1$, un camino de
    extensión (flecha discontinua) lleva al estado $F$; al alcanzar $V_1$,
    el cómputo continúa en $B_2$.}
  \label{fig:automata-conjuncion}
\end{figure}

Por ejemplo, para la fórmula $A = x_1 \wedge (x_2 \vee x_1)$, cuya
secuencia de ocurrencias es $x_1\,x_2\,x_1$, se construye primero el
autómata de cada subfórmula, $x_1$ y $x_2 \vee x_1$. El de
$x_1$ es el de una sola ocurrencia; el de $x_2 \vee x_1$ se obtiene por
la regla de la disyunción. La figura~\ref{fig:automata-partes} muestra
ambos.

\begin{figure}[ht]
  \centering
  \begin{tabular}{cc}
    \begin{tikzpicture}[->, >=stealth, shorten >=1pt, auto, semithick]
      \node[state, initial] (q0) at (0,0) {$q_0$};
      \node[state, accepting] (V) at (2.4,0.9) {$V$};
      \node[state] (F) at (2.4,-0.9) {$F$};
      \draw (q0) edge node {$1$} (V)
      (q0) edge node[swap] {$0$} (F);
    \end{tikzpicture}
                                   &
    \begin{tikzpicture}[->, >=stealth, shorten >=1pt, auto, semithick]
      \node[state, initial] (q1) at (0,0) {$q_1$};
      \node[state] (q3) at (2.4,1.1) {$q_3$};
      \node[state] (q2) at (2.4,-1.1) {$q_2$};
      \node[state, accepting] (V) at (4.8,1.1) {$V$};
      \node[state] (F) at (4.8,-1.1) {$F$};
      \draw (q1) edge node {$1$} (q3)
      (q1) edge node[swap] {$0$} (q2)
      (q3) edge node {$0,1$} (V)
      (q2) edge node {$1$} (V)
      (q2) edge node[swap] {$0$} (F);
    \end{tikzpicture}
    \\[0.4em]
    {\small (a) autómata de $x_1$} & {\small (b) autómata de $x_2 \vee x_1$}
  \end{tabular}
  \caption{Autómatas de $x_1$ y de $x_2 \vee x_1$, las dos subfórmulas de
    $A$.}
  \label{fig:automata-partes}
\end{figure}

Al combinar los dos por la regla de la conjunción se obtiene el autómata
de $A$: desde la rama verdadera de $x_1$ se entra en el autómata de
$x_2 \vee x_1$, y desde la rama falsa se llega al estado $F$ por un
camino de extensión. La figura~\ref{fig:automata-booleano} lo muestra.

\begin{figure}[ht]
  \centering
  \begin{tikzpicture}[->, >=stealth, shorten >=1pt, auto, semithick]
    \node[state, initial] (q0) at (0,0) {$q_0$};
    \node[state] (q1) at (2.6,1.3) {$q_1$};
    \node[state] (q4) at (2.6,-1.7) {$q_4$};
    \node[state] (q3) at (5.2,2.3) {$q_3$};
    \node[state] (q2) at (5.2,0.3) {$q_2$};
    \node[state] (q5) at (5.2,-1.7) {$q_5$};
    \node[state, accepting] (V) at (7.8,1.3) {$V$};
    \node[state] (F) at (7.8,-1.7) {$F$};
    \path
    (q0) edge node {$1$} (q1)
    (q0) edge node[swap] {$0$} (q4)
    (q1) edge node {$1$} (q3)
    (q1) edge node {$0$} (q2)
    (q3) edge node {$0,1$} (V)
    (q2) edge node {$1$} (V)
    (q2) edge node[swap] {$0$} (F)
    (q4) edge node {$0,1$} (q5)
    (q5) edge node {$0,1$} (F);
  \end{tikzpicture}
  \caption{Autómata booleano de la fórmula $A = x_1 \wedge (x_2 \vee
      x_1)$. El estado $V$ acepta las cadenas que hacen verdadera la fórmula
    cuando cada ocurrencia se trata por separado.}
  \label{fig:automata-booleano}
\end{figure}

Un camino del estado inicial a $V$ corresponde a una asignación que hace
verdadera la fórmula cuando cada ocurrencia se trata por separado. Por
ejemplo, con la cadena $110$ se llega del estado inicial a $V$; en ella
se asigna $1$ a la primera ocurrencia de $x_1$, $1$ a $x_2$ y $0$ a la
segunda ocurrencia de $x_1$, valores con los que $x_1 \wedge (x_2 \vee
  x_1)$ resulta verdadera. En esa misma cadena, las dos ocurrencias de
$x_1$ no toman el mismo valor, algo que no ocurre en ninguna interpretación.
El autómata booleano de una fórmula con $n$ ocurrencias de variable
\comment{tiene $O(n^2)$}{esto es verdad y es suficiente para que sea polinomial, pero cunado la formula está en forma normal conjuntiva, el autómata tiene $O(n)$ estados. Solo te digo por si querías un upper bound más bajo} estados~\cite{fernandezarias2007}. En él todas las
ocurrencias se tratan por separado, incluso las de una misma variable, y
por eso hace falta un segundo mecanismo que imponga la consistencia
entre las ocurrencias.

\agregaesto{TRANSICIÓN A LO SIGUIENTE}

\subsection{La consistencia entre ocurrencias y la clase de fórmulas}
\label{subsec:consistencia-clase}

El segundo mecanismo obliga a que las ocurrencias de una misma variable
tomen el mismo valor. Una cadena de ceros y unos, con un símbolo por
ocurrencia, es \textbf{consistente} si en las posiciones de cada variable
lleva siempre el mismo valor. La gramática genera las cadenas
consistentes de la longitud dada, y la intersección de ese lenguaje con
el del autómata booleano contiene las interpretaciones verdaderas de la
fórmula. La satisfacibilidad equivale entonces a que esa intersección no
sea vacía. SATCF y SATSRC se diferencian en el formalismo de esa
gramática, y con ello en la clase de fórmulas que resuelven.

En SATCF la consistencia se impone con una gramática libre del contexto,
\comment{equivalente a una pila}{what???? esto no es verdad, mejor di que es equivalente a un autómata de pila}~\cite{fernandezarias2007}. En la pila \agregaesto{del autómata} se guarda
el valor de una variable en su primera ocurrencia y se compara en las
siguientes. Como en una pila se accede a los valores en orden inverso al
de guardado, con este esquema la consistencia solo se verifica cuando
las ocurrencias están bien anidadas: si entre dos ocurrencias
consecutivas de una variable aparece una ocurrencia de otra, todas las
ocurrencias de esa otra deben caer entre las dos primeras. Una secuencia
de ocurrencias con esa propiedad tiene un \textbf{orden libre del
  contexto}~\cite{fernandezarias2007}. Por ejemplo, en la secuencia
$x_1\,x_2\,x_2\,x_1$ las dos ocurrencias de $x_2$ caen entre las de
$x_1$; la figura~\ref{fig:orden-ocurrencias}(a) muestra ese anidamiento.
Con SATCF se resuelven en tiempo $O(n^3)$ las fórmulas cuyo orden es de
ese tipo.

En SATSRC la consistencia se impone con una gramática de concatenación
de rango simple~\cite{aguileralopez2016}. Frente a la pila, que sigue un
solo fragmento de la cadena, una \comment{sRCG}{llegado a este punto, el lector no se acuerda de lo que es esto, sería bueno que lo recordaras en una oración o dos} admite varios a la vez, por lo que
admite cruces entre ocurrencias que el orden libre del contexto excluye.
Una secuencia para la que existe una sRCG de tamaño polinomial que
genera sus cadenas consistentes tiene un \textbf{orden de concatenación
  de rango simple}~\cite{aguileralopez2016}. Toda secuencia con orden
libre del contexto lo tiene, y además lo tienen secuencias con cruces.
Por ejemplo, en la secuencia $x_2\,x_1\,x_3\,x_2\,x_1$ las ocurrencias de
$x_2$ y las de $x_1$ se cruzan, de modo que ese orden no es libre del
contexto pero sí de concatenación de rango simple. La
figura~\ref{fig:orden-ocurrencias}(b) muestra ese cruce.
Con SATSRC se resuelven esas fórmulas en tiempo $O(n^k)$, donde $k$ es la
aridad de la sRCG, así que el grado del polinomio crece con la aridad
necesaria para el cruce.

\begin{figure}[ht]
  \centering
  \begin{tabular}{cc}
    \begin{tikzpicture}[baseline=(current bounding box.south)]
      \node (a1) at (0,0) {$x_1$};
      \node (a2) at (0.9,0) {$x_2$};
      \node (a3) at (1.8,0) {$x_2$};
      \node (a4) at (2.7,0) {$x_1$};
      \draw (a1) to[bend left=45] (a4);
      \draw (a2) to[bend left=55] (a3);
    \end{tikzpicture}
                                          &
    \begin{tikzpicture}[baseline=(current bounding box.south)]
      \node (b1) at (0,0) {$x_2$};
      \node (b2) at (0.9,0) {$x_1$};
      \node (b3) at (1.8,0) {$x_3$};
      \node (b4) at (2.7,0) {$x_2$};
      \node (b5) at (3.6,0) {$x_1$};
      \draw (b1) to[bend left=45] (b4);
      \draw (b2) to[bend left=45] (b5);
    \end{tikzpicture}
    \\[0.6em]
    {\small (a) orden libre del contexto} & {\small (b) orden con cruce}
  \end{tabular}
  \caption{Arcos que unen las ocurrencias de cada variable. En (a) los
    arcos se anidan; en (b), los de $x_2$ y los de $x_1$ se cruzan.}
  \label{fig:orden-ocurrencias}
\end{figure}

Una parte de las fórmulas ya se resuelve en tiempo polinomial. Las que
no tienen cruces entre ocurrencias tienen orden libre del contexto y se
deciden con SATCF; al aparecer cruces, la decisión sigue siendo
polinomial con SATSRC mientras la aridad necesaria esté acotada. El
problema empieza cuando no lo está: el exponente de la cota $O(n^k)$
crece con la cantidad de ocurrencias y la cota deja de ser polinomial.

Para codificar la consistencia de esas fórmulas hace falta la
no-linealidad. En una sRCG cada variable aparece una sola vez, de modo
que, para igualar dos ocurrencias, hay que seguirlas en argumentos
separados de la cláusula, y la cantidad de argumentos —la aridad— crece
con el cruce. Con la no-linealidad esa igualdad se expresa de forma
directa: una misma variable aparece más de una vez en la parte derecha
de una cláusula, y las subcadenas en que aparece deben coincidir, sin
que la aridad tenga que crecer.

La no-linealidad es justo la diferencia entre una sRCG y una RCG. Sin
ella se reconocen los mismos lenguajes que con una sRCG, así que no se
obtiene ninguna fórmula nueva. Por eso la siguiente clase natural de
esta línea son las RCG.\footnote{Entre las sRCG y las RCG está pMCFG, con una forma
  restringida de no-linealidad. Este trabajo no la aborda: ya existe un
  trabajo de la Facultad en esa dirección~\cite{artechemunoz2026}.}\comment{}{<3}

La clase es cerrada bajo intersección y el reconocimiento de la
consistencia con una RCG sigue siendo polinomial, según los
teoremas~\ref{teo:clausura-rcg} y~\ref{teo:prcg-ptime}; pero el SAT no se
decide por el reconocimiento, sino por el vacío de la intersección. En SATCF y SATSRC
esa decisión es polinomial porque el vacío de las CFG y las sRCG lo es,
no por una estrategia particular. Con la no-linealidad de las RCG la
aridad ya no crece con el cruce, pero el vacío de la intersección deja de
ser polinomial, y ahí queda concentrada la dificultad.

En este capítulo se presentaron las gramáticas de concatenación de
rango, con su reconocimiento polinomial y su clausura bajo intersección,
y el esquema por intersección con que se decide la satisfacibilidad en
SATCF y SATSRC. En el próximo capítulo se estudia ese vacío: se
construye un procedimiento explícito de decisión y se caracteriza su
costo.

\bibliographystyle{plain}
\bibliography{../bib/referencias}

\end{document}
